This appendix follows the order of the paper. Online Appendix~1 describes the invented data, the sample and the sensitivity checks behind Section~\ref{sec:data} and equation~\eqref{eq:baseline}. Tables and figures are numbered A1, A2 and so on. Equations continue the numbering of the paper. The appendix opens with a guide that answers the questions a reader is likely to ask first, each with what we did, what we found and where the evidence is.

\begin{ecofinreaderqa}
\textit{Does attrition across secondary regional rosters drive the post-break gap?} Within complete core rosters alone, the estimated post-break decline is \(-17.8\) per cent (s.e.\ \(3.1\)), compared with \(-18.4\) per cent (s.e.\ \(2.9\)) in the full panel (Table~\ref{tab:oa_checks}, row 2). The two estimates are not statistically different (\(p = 0.78\)).

\textit{Could time-varying returns to tenure explain the divergence?} When tenure profiles may vary freely across periods, as in equation~\eqref{eq:oa_decomp}, the estimate keeps \(89\) per cent of its baseline size (\(p < 0.001\); Table~\ref{tab:oa_checks}, row 3). Standard errors are clustered as in \citet{bertrand2004} and \citet{abadie2023}.
\end{ecofinreaderqa}

\ecofinappsubsection{Online Appendix 1: Sources and Specification Checks}

\paragraph{Augmented tenure specification}
To separate governance screening from time-varying tenure premia, we estimate
\begin{equation}
  y_{it} = \beta \, \mathit{Treated}_{i} \times \mathit{Post}_{t} + \theta_{t} \, \mathit{Tenure}_{it} + \alpha_{i} + \gamma_{t} + u_{it},
  \label{eq:oa_decomp}
\end{equation}
where \(\theta_{t}\) allows the slope on tenure to shift in every year. Table~\ref{tab:oa_checks} compares baseline and saturated specifications relative to Table~\ref{tab:summary}.

\begin{table}[!htbp]
\centering
\caption{Alternative Samples and Specifications}\label{tab:oa_checks}
\begin{threeparttable}
\footnotesize
\setlength{\tabcolsep}{5.5pt}
\begin{tabular}{l S[table-format=-2.1,table-space-text-post={\sym{***}}] c r}
\toprule
Specification & {Estimate (\%)} & Std.\ error & {Observations} \\
\midrule
1. Baseline full panel                      & -18.4\sym{***} & (2.9) & 144,000 \\
2. Core headquarters rosters only           & -17.8\sym{***} & (3.1) &  54,600 \\
3. Period-varying tenure slopes (eq.~\ref{eq:oa_decomp}) & -16.3\sym{***} & (3.3) & 144,000 \\
4. Office $\times$ region $\times$ year FE  & -19.1\sym{***} & (3.4) & 112,800 \\
\bottomrule
\end{tabular}
\begin{ecofintablenotes}
\item \textit{Sources:} As Table~\ref{tab:summary}.
\item \textit{Notes:} Standard errors clustered by division in parentheses. $\sym{*}p<0.10$, $\sym{**}p<0.05$, $\sym{***}p<0.01$.
\end{ecofintablenotes}
\end{threeparttable}
\end{table}

